
\section{Points de vigilance et sécurité par séance}
\label{sec:vigilance}

Cette section rassemble, \textbf{à l'usage de l'encadrant}, les points de vigilance et de
sécurité associés à chaque séance du programme. Elle est volontairement distincte du détail
pédagogique ci-dessus, qui est celui utilisé avec les élèves.

\begin{description}[leftmargin=0.7cm]
  \item[Séance 1] Adapter la profondeur du bassin ; s'assurer de l'aisance de chacun avant
    de complexifier les exercices. Vidage du tuba : rester au bord ou prévoir un
    dispositif de flottabilité, ne pas inhaler d'eau. Souffler par le nez lors des
    essais de dissociation bucco-nasale pour éviter le placage de masque.
  \item[Séance 2] Souffler par le nez pour éviter le placage de masque ; prévention du
    barotraumatisme des oreilles et des sinus. Vérifier la profondeur avant chaque saut,
    ne jamais plonger en piqué ; attention à la percussion d'un autre élève.
  \item[Séance 3] Toujours un échauffement apnée avant chaque exercice de canard ou de
    phoque, même bref. Ne jamais dépasser 5 secondes d'apnée expiratoire (phoque).
    Équilibrage des oreilles dès la surface (Valsalva en descendant, jamais en
    remontant). Tour d'horizon systématique en remontant ; attention au placage de
    masque ; travailler au bord, côté peu profond.
  \item[Séance 4] Toujours attacher le bloc ; ne jamais bloquer sa respiration ;
    expiration continue impérative pendant toute remontée (prévention de la surpression
    pulmonaire) ; vérifier l'absence d'autres plongeurs avant la mise à l'eau ; lestage
    adapté au milieu. Équilibrage des oreilles dès la surface lors de l'immersion
    (Valsalva en descendant, jamais en remontant), prévention du barotraumatisme des
    oreilles et des sinus. Purge lente : ouverture progressive pour une descente
    maîtrisée.
  \item[Séance 5] Encadrant à proximité immédiate en permanence ; respecter le rythme de
    chaque élève ; jamais de situation générant de l'insécurité. Purge haute (rapide) :
    vider progressivement.
  \item[Séance 6] Garder en permanence le détendeur en bouche pendant les exercices d'équilibrage.
    Partage d'air en statique : rester à proximité immédiate du binôme.
  \item[Séance 7] Gonflage progressif du gilet, jamais de surgonflage ; rester calme sans
    masque, ne jamais retenir sa respiration ; encadrant proche lors du retrait du
    masque.
  \item[Séance 8] Aucune recherche de performance chronométrée ; prévention de
    l'essoufflement.
  \item[Séance 9] Accent sur la prévention de la surpression pulmonaire : ne jamais
    bloquer sa ventilation pendant la remontée ; attention à l'apnée inspiratoire.
  \item[Séance 10] Exercices encadrés individuellement, procédures répétées calmement,
    jamais de simulation brutale.
  \item[Séance 11] Veiller à la cohésion de la palanquée sur l'ensemble du parcours.
  \item[Séance 12] Adapter le contenu au niveau réel de chaque élève, sans précipitation.
  \item[Séance 13] Rappel : en cas de certification délivrée en milieu artificiel, au
    moins 4 plongées en milieu naturel devront être réalisées et attestées dans les 12
    mois suivant l'obtention de la certification.
\end{description}
